\documentclass[11pt]{article}
\usepackage{amsmath,amssymb,amsthm}
\usepackage{xcolor}
\usepackage{hyperref}
\PassOptionsToPackage{hyphens}{url}
\hypersetup{colorlinks=true, linkcolor=blue!50!black, urlcolor=blue!50!black,
  citecolor=blue!50!black,
  pdftitle={The shortest relations of planar triangle reflection groups},
  pdfauthor={Vico Bonfioli}}

\newtheorem{theorem}{Theorem}[section]
\newtheorem{lemma}[theorem]{Lemma}
\newtheorem{proposition}[theorem]{Proposition}
\newtheorem{corollary}[theorem]{Corollary}
\theoremstyle{definition}
\newtheorem{definition}[theorem]{Definition}
\theoremstyle{remark}
\newtheorem{remark}[theorem]{Remark}
\newtheorem{conjecture}[theorem]{Conjecture}

\newcommand{\Z}{\mathbb{Z}}
\newcommand{\Q}{\mathbb{Q}}
\newcommand{\R}{\mathbb{R}}

\title{The shortest relations of planar triangle reflection groups}
\author{Vico Bonfioli\thanks{Independent researcher. \texttt{vicobonfioli@gmail.com}.
Code and verification data:
\url{https://github.com/ElVec1o/pythagorean-reflection-sequence}.}}
\date{\today}

\begin{document}
\maketitle

\begin{abstract}
Let $\tau$ be a Euclidean triangle and $G_\tau$ the group generated by the three
reflections in its sides. We determine the shortest relations satisfied by these
generators for every $\tau$: there are exactly $33$ of them, each equating two
reduced words of length $11$, and for all $\tau$ outside a proper Zariski-closed
set of shapes they are the only coincidences in the ball of radius $11$. The
associated growth sequence begins
$1,3,6,12,24,48,96,192,384,768,1536,3039,6012$. We complement this with the
word metric of the translation subgroup: its elements are signed
configurations of lamps on a hexagonal site lattice, and we identify them with
the finitely supported integer flows on the honeycomb Cayley graph of the
point group $\Z^2\rtimes C_2$. For all shapes outside countably many proper
subvarieties, the word length of a configuration equals the $\ell^1$-norm of
its flow plus twice the least number of edges connecting the support of the
flow to the base vertex. The cancellation and single-lamp laws follow, and the
formula reproduces the exact census of translations through length $18$,
element by element. The shape space is stratified: on the lines where one
angle equals $\pi/m$ the relation onset moves to depth $6,8,9,10,10$ for
$m=4,6,3,2,8$ and is the universal one for the remaining $m\le12$, so that
the crystallographic orders are not the distinguished ones. We prove that
each divisor $d\ge2$ of $m$ contributes $d^2-1$ relations at depth $m+m/d$,
which for even $m$ gives $(m/2)^2-1$ relations at depth $m+2$.
\end{abstract}

\medskip
\noindent\textbf{Keywords.} triangle reflection group; growth series; word
metric; Coxeter group; lamplighter group; crystallographic restriction.

\smallskip
\noindent\textbf{2020 Mathematics Subject Classification.}
Primary 20F55, 20F65; Secondary 20F05, 51F15, 05A15.

\section{Introduction}\label{sec:intro}

Let $\tau\subset\R^2$ be a nondegenerate Euclidean triangle and let
$r_0,r_1,r_2$ be the reflections in the lines containing its three sides. They
generate a group $G_\tau\le\mathrm{Isom}(\R^2)$, the \emph{triangle reflection
group} of $\tau$. When the angles of $\tau$ are of the form $\pi/n_i$ the
triangle tiles the plane and $G_\tau$ is a Euclidean reflection group with the
classical finite presentation. For all other shapes the group is not
crystallographic, and Isola and Piergallini \cite{IP} showed that for a
\emph{generic} triangle, meaning one whose edge lengths are algebraically
independent, the translation subgroup $T_\tau$ is free abelian of infinite rank,
the orientation preserving subgroup is free metabelian of rank two, and $G_\tau$
admits no finite presentation. Their work determines the relations of $G_\tau$
as a group. It does not address the word metric.

The present paper takes up the metric question. Write
$F=C_2*C_2*C_2=\langle x_0,x_1,x_2\mid x_i^2\rangle$ and let
$\rho_\tau\colon F\to G_\tau$ send $x_i$ to $r_i$. Elements of $F$ are reduced
words, and the sphere of radius $d$ in $F$ has $3\cdot2^{d-1}$ elements for
$d\ge1$. The question is where, and by how much, the image sphere is smaller.

\begin{theorem}[Census]\label{thm:census}
There is an explicit list of $33$ pairs $(w_i,w_i')$ of reduced words of length
$11$, pairwise distinct in $F$, such that
\begin{enumerate}
\item[(i)] $\rho_\tau(w_i)=\rho_\tau(w_i')$ for every nondegenerate Euclidean
triangle $\tau$ and every $i$;
\item[(ii)] for $\tau$ outside a proper Zariski-closed subset of the shape space,
these are the only coincidences of $\rho_\tau$ on the ball of radius $11$, and
on the ball of radius $12$ the only further coincidences are the $132$
described in Remark~\ref{rem:twelve}.
\end{enumerate}
Consequently the growth sequence
$u_d(\tau)=\#\{g\in G_\tau:\ \ell(g)=d\}$ of the word metric of a generic
triangle begins
\[
 1,\;3,\;6,\;12,\;24,\;48,\;96,\;192,\;384,\;768,\;1536,\;3039,\;6012,
\]
its first deviation from $3\cdot2^{d-1}$ occurring at $d=11$ with deficit $33$,
and the shortest elements of $\ker\rho_\tau$ common to all triangles have length
exactly $22$.
\end{theorem}

The second half of the paper concerns $T_\tau$. By \cite[Thm.~4.1]{IP} it is free
abelian on the conjugacy class of $t_1=(r_0r_1r_2)^2$, so its elements are
finitely supported integer vectors on a lattice of \emph{sites}; we call a
nonzero entry a \emph{lamp}. Conjugation by the generators acts on sites through
three affine involutions whose pairwise products are translations by the six
vectors of a hexagonal lattice (\S\ref{sec:sites}). In \S\ref{sec:flow} we
identify $T_\tau$, for generic shapes, with the group of finitely supported
integer flows on the Cayley graph $X$ of the point group $\Z^2\rtimes C_2$, a
cubic graph whose hexagonal faces correspond to the sites, and determine the
metric completely:

\begin{theorem}[see Theorem~\ref{thm:metric}]\label{thm:introm}
For flow-generic $\tau$ (Definition~\ref{def:generic}) and every reduced
configuration $c$,
\[
 \ell(c)\;=\;\sum_{\{u,v\}}|c_u-c_v|\;+\;2\,\mathrm{st}(c),
\]
the sum over unordered hex-adjacent pairs of sites with $c$ extended by zero,
where $\mathrm{st}(c)$ is the least number of edges of $X$ whose addition
connects the support of the flow of $c$ together with the base vertex.
\end{theorem}

The first summand is the $\ell^1$-norm of the flow of $c$; the second is a
Steiner term. The cancellation rule for lamps, the single-lamp costs of
\S\ref{sec:shield}, and the exact census of translations by word length all
follow (\S\ref{sec:flow}).

Section~\ref{sec:strata} records what happens on the non-generic strata. If one
angle is exactly $\pi/m$ and the shape is otherwise generic, then the reference
group is the Coxeter group $\langle x_i\mid x_i^2,(x_0x_1)^m\rangle$, whose
growth series we compute in closed form, and the first deficit occurs at a
depth $d^*(m)$ which, for $m\le12$, is smaller than the universal $11$ exactly
when $m\in\{2,3,4,6,8\}$. The crystallographic orders are therefore not the
ones admitting early relations, correcting the expectation that the
crystallographic restriction governs the stratification: what governs it is
the order of the rotations that the stratum creates, through the divisors of
$m$.

\section{The census theorem}\label{sec:census}

Every triangle is similar to one with vertices $(0,0)$, $(1,0)$, $(p,q)$ with
$q\neq0$, and similarity conjugates $\rho_\tau$, so all coincidence data depends
only on $(p,q)$. Over $\Q(p,q)$ the three reflections have matrix entries that
are rational functions whose denominators are the squared side lengths, and
these do not vanish for nondegenerate $\tau$.

\begin{proof}[Proof of Theorem~\ref{thm:census}]
For (i): for each of the $33$ pairs the identity $\rho(w_i)=\rho(w_i')$ is
verified by exact arithmetic in $\Q(p,q)$. An identity of rational functions
holds at every nondegenerate specialization, acute or obtuse.

For (ii): at the rational witness $\tau_0$ with third vertex $(1/3,1/2)$, exact
arithmetic over $\Q$ shows that $\rho_{\tau_0}$ is injective on the ball of
radius $10$ and that its only coincidences on the ball of radius $11$ are the
$33$ listed pairs, each consisting of two words of length exactly $11$; hence
the layer counts are $3\cdot2^{d-1}$ for $1\le d\le10$ and $3072-33=3039$ at
$d=11$. Any coincidence $\rho_\tau(w)=\rho_\tau(w')$ with $w\neq w'$ of length
at most $11$ is the vanishing of a rational function of $(p,q)$; by the witness
it is not identically zero unless $(w,w')$ is one of the $33$ pairs, so each of
the finitely many remaining conditions defines a proper closed subset. Off their
union the coincidences are exactly the $33$ pairs.

Finally, no cancellation occurs in $w_i(w_i')^{-1}$. If $w_i$ and $w_i'$ shared a
common suffix, say $w_i=us$ and $w_i'=u's$, then $\rho_\tau(u)=\rho_\tau(u')$
with $u\neq u'$ of length at most $10$, contradicting injectivity on that ball;
a common prefix is excluded in the same way. Hence each relator has length
exactly $22$, and by (ii) no shorter element of $\ker\rho_\tau$ is common to
all shapes.

For the radius $12$ statement, enumeration of all reduced words of length at
most $12$ at the witness $\tau_0$, and independently at the witness with third
vertex $(2/7,3/5)$, returns the same $132$ coincidences beyond those of the
ball of radius $11$ and no others; each of the $132$ is an identity over
$\Q(p,q)$, verified after clearing denominators as an identity of polynomials
in $\Z[p,q]$. The argument of the previous paragraph applies verbatim to the
finitely many remaining conditions on the ball of radius $12$.
\end{proof}

\begin{remark}\label{rem:twelve}
The $132$ coincidences at length $12$ carry no new relations: $66$ of them are
the one letter extensions $w_ix=w_i'x$ and $xw_i=xw_i'$ of the $33$ pairs, and
$66$ are the splittings of the $33$ relators into a word of length $12$ and a
word of length $10$, two per relator. In particular $66$ reduced words of
length $12$ represent elements at distance $10$, so the image
$\rho_\tau(\{|w|=12\})$ has $6078$ elements whereas the sphere of the word
metric has $6012$; it is the latter that the theorem records, and that the
breadth-first computations return.
\end{remark}

\begin{remark}
The witnesses $(-1/3,1/2)$, obtuse, and $(2/7,3/5)$ produce the same $33$
classes. By \cite[Thm.~4.2]{IP} each of the $33$ relators lies in the normal
closure of the commutators of $t_1$ with its conjugates, so the theorem is
consistent with, and refines metrically, the presentation given there. Their
Section~6 tabulates cyclically reduced stable words up to length $24$ modulo
symmetries; that census is of the relator language, not of the geodesic ball.
\end{remark}

\section{Sites, lamps and the hexagonal lattice}\label{sec:sites}

Normalize the incircle of $\tau$ to be the unit circle at the origin, so that
$(x)r_i=(x)s_i+2v_i$ with $s_i$ the linear part and $v_i$ the unit vector to the
tangency point of the $i$th side \cite[\S3]{IP}. Set
$t_{n,m}=(\!(t_1)\!)(r_0r_1)^n(r_0r_2)^m$; by \cite[Thm.~4.1]{IP} these form a
basis of $T_\tau$ for generic $\tau$. Conjugation by the generators acts on the
index by
\[
 \sigma_0(n,m)=(-n-1,-m+1),\quad \sigma_1(n,m)=(-n-2,-m+1),\quad
 \sigma_2(n,m)=(-n-1,-m),
\]
three affine involutions whose pairwise products are the translations by
\[
 a=(1,0),\qquad b=(0,1),\qquad c=(1,-1),\qquad a=b+c.
\]
We call $\{\pm a,\pm b,\pm c\}$ the \emph{hex generators} and say two sites are
\emph{hex-adjacent} if their difference is one of them; the sites
$(0,0),(-1,0),(-1,1)$ are the \emph{base sites}.

\begin{proposition}[Reduction]\label{prop:reduced}
Distinct reduced integer vectors on sites represent distinct elements of
$T_\tau$. In particular no site carries both signs in a reduced representative,
and a configuration containing a lamp together with its inverse at one site
represents the identity.
\end{proposition}

\begin{proof}
Immediate from the freeness of $T_\tau$ on the conjugacy class of $t_1$
\cite[Thm.~4.1]{IP}, together with injectivity of $(n,m)\mapsto t_{n,m}$ for
generic $\tau$.
\end{proof}

\section{Single lamps}\label{sec:shield}

Let $k(n,m)$ be the least length of a non-backtracking walk in the $\sigma_i$
from a base site to $(n,m)$ whose first step is by the middle letter of the
cyclic word attached to that base site.

\begin{proposition}\label{prop:shield}
$\ell(t_{n,m})\le 6+2k(n,m)$, with equality for every site with $k\le5$.
\end{proposition}

\begin{proof}
Conjugating a cyclic word $(abc)^2$ by its head or tail letter is a rotation and
preserves length; conjugating by the middle letter lengthens it by $2$, after
which the word begins and ends with that letter, so the only length preserving
conjugation is the inverse step and every further conjugation by a different
letter adds $2$. A minimal walk therefore realizes $t_{n,m}$ by a word of length
$6+2k(n,m)$. Equality on the stated range is exact breadth-first search.
\end{proof}

Corollary~\ref{cor:single} upgrades the inequality to an equality at every
site, for flow-generic shapes.

\section{The translation metric}\label{sec:flow}

Write $F=C_2*C_2*C_2$ as before and let
\[
\widehat P \;=\; \bigl\langle\, x_0,x_1,x_2 \;\bigm|\; x_i^2,\
[x_0x_1,\,x_0x_2]\,\bigr\rangle .
\]
The rotations $\rho_1=x_0x_1$ and $\rho_2=x_0x_2$ commute in $\widehat P$ and
are inverted by conjugation with $x_0$, so every element has the normal form
$\rho_1^{m_1}\rho_2^{m_2}x_0^{\varepsilon}$ and $\widehat P\cong\Z^2\rtimes
C_2$. Let $X$ be the left Cayley graph of $\widehat P$ on $x_0,x_1,x_2$, with
edges $\{g,x_ig\}$ and base vertex $e$: a connected cubic simple graph,
bipartite for the orientation character $\epsilon(x_i)=-1$. For every
nondegenerate $\tau$ the linear parts satisfy $s_i^2=1$ and
$[s_0s_1,s_0s_2]=1$, plane rotations commuting, so $x_i\mapsto s_i$ defines a
homomorphism $M_\tau\colon\widehat P\to\mathrm O(2)$ at every shape.

\subsection{Walks, flows, deposits}

For a word $w=x_{i_1}\cdots x_{i_\ell}$ let $\sigma_j\in\widehat P$ be the
image of the suffix $x_{i_{j+1}}\cdots x_{i_\ell}$, so that $\sigma_\ell=e$
and $\sigma_{j-1}=x_{i_j}\sigma_j$. The \emph{suffix walk}
$\sigma_\ell,\dots,\sigma_0$ traverses at step $j$ the edge
$\{\sigma_j,x_{i_j}\sigma_j\}$ of $X$. Orient each edge of $X$ from its
$\epsilon=+1$ endpoint and let the \emph{flow} $\phi(w)\in\Z^{E(X)}$ be the
net signed traversal count of the suffix walk. When $w$ maps to $1$ in
$\widehat P$ the walk is closed and $\phi(w)$ is a \emph{cycle}: at each
vertex the signed values of the three incident edges sum to zero. Flow is
additive under concatenation of words with trivial image.

In the incircle expansion of \S\ref{sec:sites} the translation part of
$\rho_\tau(w)$ is
\[
 t(w)\;=\;\sum_{j=1}^{\ell} 2v_{i_j}\,s_{i_{j+1}}\cdots s_{i_\ell}
 \;=\;\sum_{j=1}^{\ell} 2v_{i_j}\,M_\tau(\sigma_j).
\]
Step $j$ deposits $2v_{i_j}M_\tau(\sigma_j)$, and the deposits of the two
directions of one edge are opposite: $v_is_i=-v_i$ gives
$2v_iM_\tau(x_i\sigma)=-2v_iM_\tau(\sigma)$. Writing $d_\tau(E)$ for the
deposit of the canonical direction of the edge $E$,
\begin{equation}\label{eq:flow}
 t(w)\;=\;V_\tau\bigl(\phi(w)\bigr),\qquad
 V_\tau(\phi)\;=\;\sum_{E\in E(X)}\phi_E\, d_\tau(E),
\end{equation}
for every shape and every word with trivial image.

\subsection{Faces}

Let $w_6=(x_0x_1x_2)^2$, the word of $t_1$.

\begin{lemma}\label{lem:kernel}
Modulo the relations $x_i^2$ one has $w_6=[\rho_1,\rho_2^{-1}]$; consequently
$\ker(F\to\widehat P)$ is the normal closure of $w_6$, and the flows of its
elements form the $\Z$-span of the flows of the conjugates $uw_6^{\pm1}u^{-1}$.
\end{lemma}

\begin{proof}
$x_0x_1x_2x_0x_1x_2=(x_0x_1)(x_2x_0)(x_1x_0)(x_0x_2)
=\rho_1\rho_2^{-1}\rho_1^{-1}\rho_2$. Hence imposing $w_6=1$ on the free
product imposes exactly the commutation of $\rho_1$ and $\rho_2$, which
presents $\widehat P$. A product of conjugates of $w_6^{\pm1}$ is a
concatenation of words with trivial image, on which flow is additive.
\end{proof}

\begin{lemma}[Faces]\label{lem:faces}
The support of $\phi(w_6)$ is a hexagonal cycle through $e$, and
$\phi(uw_6^{\pm1}u^{-1})=\pm\,\phi(w_6)\cdot g^{-1}$, where $g$ is the image
of $u$ and $\widehat P$ acts on $X$ by right translation. The translated
hexagons, the \emph{faces} of $X$, are indexed by the sites of
\S\ref{sec:sites} through the dictionary $t_{n,m}$; write $\partial f_s$ for
the oriented flow of the face at site $s$. Then:
\begin{enumerate}
\item[(i)] every vertex of $X$ lies on exactly three faces, the faces at $e$
being those of the three base sites, and every edge lies on exactly two;
\item[(ii)] two faces share an edge exactly when their sites are
hex-adjacent, and then they share one edge, with opposite orientations;
\item[(iii)] the $\partial f_s$ are linearly independent, and the flows of
the elements of $\ker(F\to\widehat P)$ are precisely the finite sums
$\sum_s c_s\,\partial f_s$ with $c_s\in\Z$.
\end{enumerate}
\end{lemma}

\begin{proof}
In the suffix walk of $uw_6^{\pm1}u^{-1}$ the suffixes of the $u^{-1}$
segment retrace those of the $u$ segment in the opposite direction, so their
traversals cancel, while the suffixes of the middle segment are those of
$w_6^{\pm1}$ multiplied on the right by $g^{-1}$; this is the displayed
formula. The suffix walk of $w_6$ visits six distinct elements of $\widehat
P$ and closes, giving a hexagon through $e$. Conjugation permutes the lamps
through the involutions $\sigma_i$ of \S\ref{sec:sites}, which indexes the
orbit of hexagons by the sites. Right translation is a graph automorphism
acting transitively on vertices, so (i) follows from its validity at $e$,
where exact computation shows the three faces through $e$ are those of the
base sites $(0,0),(-1,0),(-1,1)$ and each of the three edges at $e$ lies on
two of them. For (ii), faces sharing an edge have intersecting vertex sets,
which confines the difference of their sites to a finite set; exact
computation over this set shows a common edge occurs precisely at the six hex
differences, is then unique, and carries opposite orientations in the two
faces. For (iii), spanning is Lemma~\ref{lem:kernel}. Independence: order
sites lexicographically and let $s_0$ be maximal in the support of
$\sum_s c_s\partial f_s$; the edge shared by $f_{s_0}$ and $f_{s_0+a}$
receives $\pm c_{s_0}\neq0$, since $c_{s_0+a}=0$.
\end{proof}

\subsection{Separation and genericity}

\begin{lemma}[Separation]\label{lem:sep}
For every nonzero finitely supported cycle $\phi$ on $X$, the set of shapes
$(p,q)$ with $V_\tau(\phi)=0$ is a proper Zariski-closed subset of the shape
plane.
\end{lemma}

\begin{proof}
The entries of $V_\tau(\phi)$ are rational functions of $(p,q)$, so it
suffices to show that $V_\tau(\phi)$ does not vanish identically. Every edge
$E$ has one endpoint with $\epsilon=+1$; write its normal form
$\rho_1^{m_1}\rho_2^{m_2}$ and let $i$ be the letter of $E$. Let
$\beta_1,\beta_2$ be the rotation angles of $M_\tau(\rho_1),M_\tau(\rho_2)$
and $\theta_i$ the direction of the $i$th side, so that
$\theta_1\equiv\theta_0-\beta_1/2$ and $\theta_2\equiv\theta_0-\beta_2/2$
modulo $\pi$, up to one global choice of sign of the $\beta_k$. The deposit
$d_\tau(E)$ is a nonzero vector in the direction
$\theta_i+\tfrac\pi2+\eta\,(m_1\beta_1+m_2\beta_2)$ modulo $\pi$, with
$\eta=\pm1$ one global sign. The \emph{frequency} of $E$, the coefficient
vector of $(\beta_1,\beta_2)$ in this direction, lies in
$(\tfrac12\Z)^2$; doubled, the frequencies of edges with letter $0,1,2$ are
congruent to $(0,0),(1,0),(0,1)$ modulo $2$ respectively. With the two global
signs fixed, the letter of $E$ is determined by the parity of its doubled
frequency and the exponents $(m_1,m_2)$ by its integer part, so distinct
edges have distinct frequencies. If $V_\tau(\phi)$ vanished identically then, the map from
shapes to $(\beta_1,\beta_2)$ being open near a suitable shape, a finite sum
$\sum_E \pm\phi_E\,|d_\tau(E)|\,e^{i(c_E+f_E\cdot\beta)}$ with distinct
frequencies $f_E$ would vanish on a nonempty open set of $\beta$;
substituting $z_k=e^{i\beta_k/2}$ makes it a Laurent polynomial vanishing on
an open subset of the torus, so all coefficients vanish, and
$|d_\tau(E)|\neq0$ forces $\phi=0$.
\end{proof}

\begin{definition}\label{def:generic}
The shape $\tau$ is \emph{flow-generic} if $M_\tau$ is injective and
$V_\tau(\phi)\neq0$ for every nonzero finitely supported cycle $\phi$ on $X$.
\end{definition}

\begin{proposition}\label{prop:generic}
Flow-genericity fails only on a countable union of proper Zariski-closed
subsets of the shape plane. In particular every $\tau$ whose coordinates
$p,q$ are algebraically independent over $\Q$ is flow-generic.
\end{proposition}

\begin{proof}
Injectivity of $M_\tau$ reduces to injectivity on the rotation subgroup,
reflections never being trivial, and there it is the condition
$m_1\beta_1+m_2\beta_2\notin2\pi\Z$ for $(m_1,m_2)\neq(0,0)$; each failure
is a nontrivial algebraic condition holding at some shapes and not at
others, hence a proper Zariski-closed set, and there are countably many.
With Lemma~\ref{lem:sep} this exhausts the conditions, each defined over
$\Q$; a point with algebraically independent coordinates lies on no proper
$\Q$-defined Zariski-closed subset.
\end{proof}

\begin{theorem}[Structure]\label{thm:structure}
Let $\tau$ be flow-generic. Every word representing a translation has closed
suffix walk; all words representing the same translation have the same flow;
and $c\mapsto\phi(c):=\sum_s c_s\,\partial f_s$ identifies the finitely
supported integer vectors on sites both with the cycles of
Lemma~\ref{lem:faces}(iii) and with $T_\tau$. In particular $T_\tau$ is free
abelian on the lamps, recovering \cite[Thm.~4.1]{IP} for flow-generic
shapes, and
\begin{equation}\label{eq:norm}
 \|\phi(c)\|_1\;=\;\sum_{\{u,v\}}|c_u-c_v|,
\end{equation}
the sum over unordered hex-adjacent pairs of sites with $c$ extended by zero.
\end{theorem}

\begin{proof}
A word $w$ with $\rho_\tau(w)\in T_\tau$ has trivial linear part, hence
trivial image in $\widehat P$ by injectivity of $M_\tau$, and $\phi(w)$ lies
in the span of the faces by Lemma~\ref{lem:faces}(iii). If $w,w'$ represent
the same translation then $V_\tau(\phi(w)-\phi(w'))=0$ by~\eqref{eq:flow},
so $\phi(w)=\phi(w')$ by flow-genericity. Taking products of conjugated lamp
words matches the flow of the configuration $c$ with $\sum_s c_s\partial
f_s$, injectivity of $c\mapsto\phi(c)$ is Lemma~\ref{lem:faces}(iii), and
surjectivity onto $T_\tau$ follows since every translation is represented by
some word. For~\eqref{eq:norm}: every edge of the face of $u$ is shared with
exactly one neighbouring face $v$ by Lemma~\ref{lem:faces}(i), and carries
the flow $\pm(c_u-c_v)$ by (ii).
\end{proof}

\subsection{The metric theorem}

\begin{theorem}[Metric]\label{thm:metric}
Let $c$ be a nonzero reduced configuration with flow $\phi=\phi(c)$, and let
$\mathrm{st}(c)$ be the least cardinality of a set $M$ of edges of $X$
outside the support of $\phi$ such that the subgraph of $X$ consisting of
the edges $\operatorname{supp}\phi\cup M$, their endpoints, and the vertex
$e$ is connected. Then
\[
 \ell(c)\;=\;\|\phi(c)\|_1\;+\;2\,\mathrm{st}(c)
\]
for flow-generic $\tau$, and $\ell(c)\le\|\phi(c)\|_1+2\,\mathrm{st}(c)$ for
every $\tau$. In particular $\ell(c)\equiv\|\phi(c)\|_1\pmod2$, and every
translation has even length.
\end{theorem}

\begin{proof}
Lower bound. Let $w$ be any word representing $c$ and let $n_E$ count all
traversals of the edge $E$ by its suffix walk, so that $\ell(w)=\sum_E n_E$,
and, the two directions of an edge contributing opposite signs,
$n_E\ge|\phi_E|$ and $n_E\equiv\phi_E\pmod2$, where $\phi=\phi(w)=\phi(c)$
by Theorem~\ref{thm:structure}. The traversed edges form a connected
subgraph containing $e$, consecutive steps sharing a vertex and the walk
beginning at $e$. Let $M$ be the set of traversed edges with $\phi_E=0$;
each has $n_E\ge2$ by parity, the set
$\operatorname{supp}\phi\cup M\cup\{e\}$ is connected, so
$|M|\ge\mathrm{st}(c)$, and
\[
 \ell(w)\;=\;\sum_{\operatorname{supp}\phi}n_E+\sum_{M}n_E
 \;\ge\;\|\phi\|_1+2|M|\;\ge\;\|\phi\|_1+2\,\mathrm{st}(c).
\]
The congruence follows by summing $n_E\equiv\phi_E\pmod2$, and evenness
from~\eqref{eq:norm}, each site having six incident lattice edges.

Upper bound. Let $M$ attain $\mathrm{st}(c)$. Form the directed multigraph
with $|\phi_E|$ parallel copies of each support edge, directed by the sign
of $\phi_E$, and one copy of each direction of each edge of $M$. Every
vertex is balanced: the cycle condition at a vertex is the balance of the
signed values of its incident edges, and the doubled edges of $M$ are
balanced. The multigraph is connected, contains $e$, and $e$ is incident to
an edge, so it carries an Eulerian circuit through $e$; reading the letters
of its edges in reverse order yields a word whose suffix walk is the
circuit, of length $\|\phi\|_1+2|M|$, with flow $\phi$ and trivial image in
$\widehat P$. Its linear part is trivial and its translation part is
$V_\tau(\phi(c))$ by~\eqref{eq:flow}, so it represents $c$, at every shape.
\end{proof}

\begin{corollary}[Single lamps]\label{cor:single}
For flow-generic $\tau$,
$\ell(t_{n,m}^{\pm1})=6+2\,d_X(e,f_{n,m})$, the graph distance in $X$ from
the base vertex to the vertex set of the face of $(n,m)$. The values
$d_X(e,f_{n,m})$ agree with the walk count $k(n,m)$ of \S\ref{sec:shield}
throughout the computed range $|n|,|m|\le7$.
\end{corollary}

\begin{proof}
Here $\|\phi\|_1=6$. A geodesic from $e$ to the nearest vertex of the face
uses no face edge and connects, so $\mathrm{st}\le d_X(e,f_{n,m})$;
conversely the connected graph $\operatorname{supp}\phi\cup M\cup\{e\}$
contains a path from $e$ to the face whose edges, up to its first face
vertex, all lie in $M$, so $\mathrm{st}\ge d_X(e,f_{n,m})$. The agreement with $k$ is exact
computation.
\end{proof}

\begin{corollary}[Cancellation]\label{cor:cancel}
$\|\phi(c)\|_1=6\sum_s|c_s|-2\sum_{\{u,v\}}\min(|c_u|,|c_v|)$, the second
sum over hex-adjacent pairs of sites carrying equal signs. For
configurations without repeated lamps this is $6L-2h$, with $L$ the number
of lamps and $h$ the number of same-sign hex-adjacent site pairs.
\end{corollary}

\begin{proof}
On the edge between adjacent sites $u,v$ one has $|c_u-c_v|=|c_u|+|c_v|$
when the signs differ or one side vanishes, and
$|c_u|+|c_v|-2\min(|c_u|,|c_v|)$ when they agree; sum over the six edges of
each site and halve the double count.
\end{proof}

\begin{corollary}[Connected stratum]\label{cor:connected}
$\mathrm{st}(c)=0$ if and only if the support of $\phi(c)$ together with $e$
is connected, and then $\ell(c)=\sum_{\{u,v\}}|c_u-c_v|$.
\end{corollary}

Two families illustrate the connectivity hypothesis. For the configuration
with one lamp of a common sign at each of the three base sites, the three
edges at $e$ are the edges shared by the three base faces and their flows
cancel, so $e$ is isolated from the support: $\mathrm{st}=1$ and
$\ell=12+2=14$. For the twelve sites at hex-distance two from a fixed site,
each carrying one lamp of the same sign, the support of the flow consists of
two vertex-disjoint cycles, so $\mathrm{st}\ge1$ for every translate of the
configuration, although the site set is hex-connected; translated to contain
a base site it has $\mathrm{st}=1$ and $\ell=48+2=50$. Connectivity of the
site set therefore does not suffice, and the support condition of
Corollary~\ref{cor:connected} cannot be weakened.

The number of translations of each word length for flow-generic $\tau$ is,
as computed below,
\[
 6,\;6,\;42,\;96,\;350,\;1092,\;3684
\]
at lengths $6,8,10,12,14,16,18$ respectively, and odd lengths carry none, as
the theorem requires: enumeration of the configurations with
$\|\phi(c)\|_1+2\,\mathrm{st}(c)\le18$, over the bounded region described in
the repository script, yields these counts, and exact rational breadth-first
search at the witness with apex $(2/7,3/5)$ returns the same sets of
translation vectors, element by element at every length, with distinct
configurations receiving distinct vectors. At the witness with
apex $(1/3,1/2)$ the same search returns $3680$ at length $18$: that shape is
not flow-generic, and four coincidences mask the count there. The
configurations with one positive lamp at each of $(0,-2),(0,-1),(1,-2),(2,-2)$
and at each of $(0,1),(1,1),(2,0),(2,1)$, of length $18$ at flow-generic
shapes, represent the same translation at that witness, as do the
configurations on $(-2,-1),(-2,0),(-1,-1),(0,-1)$ and on
$(-2,2),(-1,2),(0,1),(0,2)$, together with the inverses of both pairs; at
every other length the two witnesses agree. Of the $500$ translations of
length at most $14$, the values $\mathrm{st}=0,1,2,3,4$ occur
$264,152,54,12,18$ times.

\section{Non-generic strata}\label{sec:strata}

\begin{proposition}\label{prop:coxeter}
The growth series of $W_m=\langle x_0,x_1,x_2\mid x_i^2,(x_0x_1)^m\rangle$ with
respect to $x_0,x_1,x_2$ is
\[
  W_m(t)=\frac{(1+t)\,(1+t+\cdots+t^{m-1})}{1-t-t^2-\cdots-t^m},
\]
so its growth rate is the $m$-bonacci constant.
\end{proposition}

\begin{proof}
The parabolic formula for a Coxeter system with one finite rank two parabolic of
order $2m$ and two infinite bonds gives
$1/W_m(t^{-1})=1-3/(1+t)+1/\bigl((1+t)(1+t+\cdots+t^{m-1})\bigr)$; solving for
$W_m$ yields the displayed expression.
\end{proof}

Let one angle of $\tau$ equal $\pi/m$, the shape being otherwise generic, and
let $d^*(m)$ be the first depth at which $u_d$ falls below $[t^d]W_m(t)$, with
deficit $\delta(m)$ there. Exact computation in $\Q(2\cos(\pi/m))$, at two
rational choices of the legs for each $m$, gives
\[
\begin{array}{r|rrrrrrrrrrr}
 m & 2 & 3 & 4 & 5 & 6 & 7 & 8 & 9 & 10 & 11 & 12\\\hline
 d^* & 10 & 9 & 6 & 11 & 8 & 11 & 10 & 11 & 11 & 11 & 11\\
 \delta & 8 & 3 & 3 & 33 & 8 & 33 & 15 & 33 & 33 & 33 & 33
\end{array}
\]
The entries $(11,33)$ are the universal deficit of Theorem~\ref{thm:census}.
Beyond the universal values, the deficit at depth $12$ exceeds the generic
$132$ by $6$ for $m=5$, by $8$ for $m=9$ and by $24$ for $m=10$, and equals it
for $m=7,11,12$.

The onset is therefore earlier than the universal depth $11$ exactly for
$m\in\{2,3,4,6,8\}$ in this range. The crystallographic orders are not
distinguished: $m=8$ has onset $10$, and the mechanism is instead one of
depth, as the following theorem shows.

\begin{theorem}[Rotation relations]\label{thm:rot}
Let $m\ge3$ and let $\tau$ have angle $\pi/m$ at the vertex where the sides $0$
and $1$ meet, the shape being otherwise arbitrary. For a word $w$ in
$r_0,r_1,r_2$ and a letter $i$, let $c_i(w)$ be the number of occurrences of
$i$ at odd positions of $w$ minus the number at even positions. Fix
$1\le c\le m-1$ and put $n=m/\gcd(c,m)$.
\begin{enumerate}
\item[(i)] Let $w$ have even length. Its linear part is, at every shape of the
stratum, the rotation by $2\pi c_1(w)/m$ if and only if $c_2(w)=0$; such a $w$ represents a
rotation of order $m/\gcd(c_1(w),m)$ if $m\nmid c_1(w)$, and a translation
otherwise.
\item[(ii)] The reduced words of length $2c+2$ with $c_2=0$ and $c_1=c$ are
exactly the $c^2$ words $w_{a,b}$ with $0\le a\le c$, $1\le b\le c+1$ and
$b\notin\{a,a+1\}$, in which every odd position carries $1$ except position
$2a+1$, carrying $2$, and every even position carries $0$ except position
$2b$, carrying $2$. Reversal maps them bijectively onto the reduced words of
length $2c+2$ with $c_2=0$ and $c_1=-c$, and no shorter word of even length has $c_2=0$
and $|c_1|=c$ except $(r_0r_1)^c$ and its reversal.
\item[(iii)] Exactly one of the $w_{a,b}$, namely $w_{0,c+1}=r_2(r_0r_1)^cr_2$,
has finite order in $W_m$. Each of the other $c^2-1$ satisfies $w^n=1$
throughout the stratum without this being a consequence of the Coxeter
relations, and $w^n$ is cyclically reduced of length $n(2c+2)$, so splitting it in half
exhibits a coincidence between two reduced words of length $n(c+1)$; they are
distinct, since equality in $F$ would give $w^n=1$ there. Whether the
coincidence lowers the sphere at that depth requires the two halves to be
geodesic in $W_m$; this holds in the range $c\le m/2$ used in
Corollary~\ref{cor:onset}, where every $\{0,1\}$-block of $w_{a,b}$, including
the join $u_2u_0$, has length at most $m$, and can fail for larger $c$.
\end{enumerate}
\end{theorem}

\begin{corollary}\label{cor:onset}
Let $d$ be a divisor of $m$ with $2\le d\le m/2$. The stratum carries $d^2-1$
relations at depth $m+m/d$. Taking $d$ largest, $d^{*}(m)\le m+p$ with $p$ the
least prime factor of $m$, the family having $(m/p)^2-1$ members. For $m$
prime no such $d$ exists, and minimising the depth $(c+1)m/\gcd(c,m)$ of
Theorem~\ref{thm:rot} over $2\le c\le m-1$ instead gives $c=2$, so
$d^{*}(m)\le3m$ with $3$ members. For even $m$ this is
$d^{*}(m)\le m+2$ with $(m/2)^2-1$ members, and the relations are half-turns,
$n=2$ making each relation $w=\overline{w}$.
\end{corollary}

\begin{proof}
(i) For even $L$ the linear part of $r_{a_1}\cdots r_{a_L}$ is the rotation by
$2\sum_j(-1)^{j+1}\theta_{a_j}=2\bigl(c_0\theta_0+c_1\theta_1+c_2\theta_2\bigr)$,
with $\theta_i$ the direction of the $i$th side; normalize $\theta_0=0$ and
$\theta_1=\pi/m$. As the shape runs through the stratum $\theta_2$ runs through
an interval while $\theta_0,\theta_1$ stay fixed, so the linear part is the
same at every shape of the stratum precisely when $c_2=0$, and then it is the
rotation by $2\pi c_1/m$. A plane isometry whose linear part is a nontrivial
rotation fixes a point and is the rotation by that angle about it, whence the
stated order; if the linear part is trivial the isometry is a translation.

(ii) Put $L=2c+2$ and let $p,q$ count the letters $1$ at odd and at even
positions, so $p-q=c_1=c$. Each parity class has $c+1$ positions, so $p\le c+1$
and hence $q\le1$. If $q=1$ then $p=c+1$, so every odd position carries $1$;
the even position carrying $1$ has a neighbour, necessarily odd and carrying
$1$, contradicting reducedness. So $q=0$ and $p=c$: no even position carries
$1$, and exactly one odd position carries another letter. Were that letter $0$,
then no letter $2$ would occur at an odd position, hence none at an even
position either since $c_2=0$, so every even position would carry $0$ and the
exceptional odd position would have a neighbour carrying $0$, again
contradicting reducedness. Hence it carries $2$, and $c_2=0$ forces exactly one
even position to carry $2$, the others carrying $0$. This is the displayed
shape. Letters $2$ at positions $2a+1$ and $2b$ are adjacent exactly when
$b\in\{a,a+1\}$, and no other neighbouring letters coincide, so the admissible
pairs are those listed, $(c+1)^2-(2c+1)=c^2$ in number. Reversal of a word of
even length exchanges the two parity classes and so negates every $c_i$. For
the last clause, $|c_1|=c$ requires at least $c$ letters $1$ in one parity
class, so an even length is at least $2c$; at length $2c$ the argument just
given leaves only $(r_0r_1)^c$ and its reversal. (For odd lengths the linear
part is a reflection and (i) does not apply.)

(iii) By (i) each $w_{a,b}$ is a rotation of order $n$ at every shape of the
stratum, so $w_{a,b}^n=1$ there. Since the Coxeter presentation of $W_m$
imposes no relation involving $x_2$ beyond $x_2^2$,
\[
 W_m\;\cong\;D_m*C_2,\qquad D_m=\langle x_0,x_1\rangle,\quad C_2=\langle
 x_2\rangle,
\]
and by the torsion theorem for free products every element of finite order is
conjugate into a factor; equivalently a cyclically reduced element with two or
more syllables from $D_m$ has infinite order. Write
$w_{a,b}=u_0x_2u_1x_2u_2$ with $u_i\in D_m$ the images of its maximal
$\{0,1\}$-blocks. Its cyclic reduction is $u_2u_0\,x_2u_1x_2$, so, $u_1$ being
nontrivial by reducedness, $w_{a,b}$ has finite order if and only if
$u_2=u_0^{-1}$. Suppose first $b\ge a+2$. Reading the blocks gives
$u_0=(x_1x_0)^a$, $u_1=(x_0x_1)^{b-a-1}$ and $u_2=(x_1x_0)^{c+1-b}$, so
$u_2=u_0^{-1}$ amounts to $b\equiv a+c+1\pmod{m}$; as $2\le b-a\le c+1\le m$
this forces $b-a=c+1$, that is $a=0$ and $b=c+1$, where indeed
$w_{0,c+1}=r_2(r_0r_1)^cr_2$ is conjugate to $(r_0r_1)^c$. Suppose next
$b\le a-1$. Then $u_0=(x_1x_0)^{b-1}x_1$, $u_1=(x_1x_0)^{a-b}$ and
$u_2=x_0(x_1x_0)^{c-a}$; writing $s=x_0x_1$ and $x_1=x_0s$, the condition
$u_2=u_0^{-1}$ becomes $x_0s^{a-c}=x_0s^{b}$, that is $b\equiv a-c\pmod{m}$,
incompatible with $1\le b\le a-1\le c-1$. Finally the first letter of
$w_{a,b}$ is $1$ unless $a=0$ and the last is $0$ unless $b=c+1$, so for the
remaining $c^2-1$ words no cancellation occurs in $w_{a,b}^n$, which is
therefore cyclically reduced of length $n(2c+2)$.
\end{proof}

The counts $3,8,15,24$ at depths $6,8,10,12$ for $m=4,6,8,10$ agree with
Corollary~\ref{cor:onset}; $m=2$ gives none, and $m=12$ is placed at depth
$14$ with $35$ relations, beyond the computed window and consistent with the
absence of extra relations at depth $12$. For odd $m$ the corollary gives $3$
relations at depth $9$ for $m=3$ and $8$ at depth $12$ for $m=9$, both
attained, and for the primes $m=5,7,11$ it gives $3$ relations at depths
$15,21,33$. The bound is attained for every computed onset except two. For
$m=2$ the theorem produces nothing, and for $m=5$ six relations occur already
at depth $12$, before the predicted $15$. Of those six, three equate two
translations: on the stratum $(r_0r_1)^m=1$ identifies $t_{n,j}$ with
$t_{n+m,j}$, so the site lattice of \S\ref{sec:sites} wraps into the cylinder
$(\Z/m)\times\Z$ and the flows of \S\ref{sec:flow} live on a cylindrical
honeycomb; the remaining three have $c_2=\pm1$. Theorem~\ref{thm:cyl} below
accounts for the translation kind: they are pairs of words whose flows differ
by a wrapping cycle. We do not count them here, and for $m=7$ and $m=11$ we do
not know whether relations of this second kind precede the depths $21$ and
$33$ of the theorem. That the relations of
Theorem~\ref{thm:rot} are the only coincidences at their depth, so that
$\delta$ equals the stated count rather than merely dominating it, is the
computational input recorded above.

For $m=3,4,6,8$ the relations tabulated above hold on the whole stratum: each
is verified identically in the legs $(a,b)$ over the relevant field, so the
onset is a property of the stratum and not of the samples, while the deficit
being exactly $\delta$ is witnessed at the rational samples and therefore
holds off a proper closed subset of the stratum. The row $m=2$ is verified at
three rational samples. The three relations at depth $9$ for $m=3$ are
\[
 210201210=102012102,\quad 020121020=121020121,\quad 201210201=012102012,
\]
of which the first and third have the form $R\,s\,R\,s^{-1}R\,s$ with $R=r_2$
and $s=r_1r_0$ the apex rotation of order three, while the second has as
relator the cube $(r_0r_2r_0r_1r_2r_1)^3$. The three at depth $6$ for $m=4$
are
\[
 012102=201210,\quad 020121=121020,\quad 102012=210201 .
\]

\subsection{The translation subgroup on a stratum}

Let $\tau$ be generic on the stratum $\alpha=\pi/m$ and put
$\widehat P_m=\widehat P/\langle\rho_1^m\rangle\cong(\Z/m\times\Z)\rtimes C_2$,
with Cayley graph $X_m$: the quotient of the honeycomb $X$ by the $m$-fold
translation, a cylinder. Its faces are the wrapped sites $(n\bmod m,\,j)$, and
the constructions of \S\ref{sec:flow} apply to $X_m$ unchanged, a word
representing a translation having a closed suffix walk whose flow is a finitely
supported cycle on $X_m$ with $t(w)=V_\tau(\phi(w))$.

\begin{definition}\label{def:stratgen}
Call $\tau$ \emph{stratum-generic} for $\alpha=\pi/m$ if $M_\tau$ is injective
on $\widehat P_m$. Since $M_\tau$ is injective on the rotation subgroup exactly
when $a\,(2\pi/m)+b\,\beta_2\notin2\pi\Z$ for $(a\bmod m,b)\neq(0,0)$, and
each failure is a nontrivial condition on the remaining shape parameter, this
holds off countably many proper closed subsets of the stratum.
\end{definition}

\begin{theorem}[Stratum metric]\label{thm:cyl}
Let $\tau$ be stratum-generic for $\alpha=\pi/m$.
\begin{enumerate}
\item[(i)] $\sum_{n\in\Z/m}t_{n,j}=0$ for every $j$: the $m$ lamps of a level
sum to zero, so the wrapped lamps are not independent.
\item[(ii)] The flow of a word representing a given translation is no longer
unique. The word $(r_0r_1)^m$ and its conjugates have closed suffix walks,
whose flows wrap once around the cylinder and are killed by $V_\tau$. Write
$\gamma_j$ for the flow of $(\!(r_0r_1)^m\!)(r_0r_2)^j$, the conjugate by the
$j$th power of $r_0r_2$. Then $\gamma_j-\gamma_{j-1}$ is the sum of the $m$
faces of level $j$. Here $Z_1(X_m)$ denotes the finitely supported cycles.
\item[(iii)] For every translation $g$,
\[
 \ell(g)\;=\;\min\bigl\{\,\|\phi\|_1+2\,\mathrm{st}(\phi)\;:\;\phi\in Z_1(X_m),\
 V_\tau(\phi)=g\,\bigr\}.
\]
\end{enumerate}
\end{theorem}

\begin{proof}
(i) Conjugation by $r_0r_1$ is the rotation by $2\pi/m$ about the apex, so it
rotates the translation vector of a lamp by that angle; the $m$ vectors
$t_{n,j}$ with $n\in\Z/m$ are therefore the images of one vector under the
$m$th roots of unity, and those sum to zero. (ii) $(r_0r_1)^m$ is the rotation
by $2\pi$, that is the identity, so the flow of its word has $V_\tau=0$; the
identity between $\gamma_j-\gamma_{j-1}$ and the level sum holds because both
are checked to agree as flows, a finite computation recorded in the script
cited below.
(iii) Neither bound of Theorem~\ref{thm:metric} used the planarity of $X$ or
the uniqueness of the flow. The lower bound gives
$\ell(w)\ge\|\phi(w)\|_1+2\,\mathrm{st}(\phi(w))$ for every word $w$
representing $g$, and $\phi(w)$ is one of the competitors; conversely the
Eulerian construction turns any competitor $\phi$ into a word of length
$\|\phi\|_1+2\,\mathrm{st}(\phi)$ with flow $\phi$, representing
$V_\tau(\phi)=g$.
\end{proof}

For $m=5$ the formula was checked against exact rational breadth-first search
on all $158$ translations of length at most $12$, the minimum being attained
already among the flows differing from $\phi(w)$ by a single $\gamma_j$. That
stratum has $6,6,42,104$ translations at lengths $6,8,10,12$, against
$6,6,42,96$ for a generic shape: wrapping shortens elements as well as
identifying them.

The identifications themselves can be counted. Sending each of the $5276$
generic translations of length at most $18$ to its image on the stratum, the
first coincidence appears at generic length $2m+4$ for $m=3,5,7$, that is at
$10,14,18$. For $m=5$ and $m=7$ it is the plain wrap, the sites $(2,0)$ and
$(-3,0)$ for $m=5$, and $(3,0)$ and $(-4,0)$ for $m=7$, which are identified
by the cylinder and are the first such pair of equal generic length; for $m=3$
it is instead the level relation of part (i), which equates the lamp $t_{0,0}$
of length $6$ with the configuration $-t_{-2,0}-t_{-1,0}$ of length $10$.
Within that window the numbers of identifications are $3255$ for $m=3$, $476$
for $m=5$, $14$ for $m=7$, and none for $m=9$ and $m=11$; for $m=3$ moreover
$76$ generic translations become trivial.

These identifications are of two kinds. Those between configurations with the
same wrapped configuration hold already in $W_m$, where $(x_0x_1)^m=1$ forces
$t_{n,j}=t_{n+m,j}$; they are $1438$, $66$ and $2$ for $m=3,5,7$ and are not
new relations. The remaining $1742$, $410$ and $12$ come from the level
relation of Theorem~\ref{thm:cyl}(i) and are new. The first of them is in each
case one level split into two arcs,
\[
 t_{0,0}=-t_{1,0}-t_{2,0}\ (m=3),\quad
 t_{3,0}+t_{4,0}=-\!\!\sum_{n\le2}\!t_{n,0}\ (m=5),\quad
 \sum_{n\le2}t_{n,0}=-\!\!\sum_{n\ge3}\!t_{n,0}\ (m=7),
\]
at generic lengths $6$ and $10$, $10$ and $14$, $14$ and $18$ respectively. An
arc of $k$ consecutive lamps of one sign has length $4k+2$, by
Corollary~\ref{cor:cancel} and Corollary~\ref{cor:connected}, so the balanced
split of a level of $m$ lamps first becomes visible at depth
$4\lceil m/2\rceil+2$, which is $2m+4$ for odd $m$ and $2m+2$ for even $m$.

Comparing with Corollary~\ref{cor:onset}: for even $m$ the rotation relations
at depth $m+2$ come first, and for odd composite $m$ so does the divisor bound
$m+p$; for $m=3$ the rotation depth $9$ precedes $10$. But for every prime
$m\ge5$ one has $2m+4<3m$, so there the level relations arrive first, at
depths $14,18,26$ for $m=5,7,11$. The stratum also carries translations that
are not images of generic translations, a word being able to have trivial
linear part on the stratum without having one generically. Both effects are
accounted for by the following census.

\begin{lemma}[Walk counts on the base rows]\label{lem:krows}
$k(n,0)=2n$ for $n\ge0$ and $-2n-2$ for $n\le-1$; $k(n,1)=2n+1$ for $n\ge0$,
$k(-1,1)=0$, and $k(n,1)=-2n-3$ for $n\le-2$. On every row $j$ the function
$n\mapsto k(n,j)$ increases by $2$ with each unit step away from its minimum,
and that minimum is $0$ for $j\in\{0,1\}$ and at least $1$ otherwise.
\end{lemma}

\begin{proof}
The three involutions $\sigma_i$ each reverse the sign of $n$ and shift it, so
a non-backtracking walk alternates the two orientations and moves $|n|$ by one
per two steps, changing $j$ only through $\sigma_0,\sigma_1$; the displayed
values are the resulting distances from the base sites $(0,0)$, $(-1,0)$,
$(-1,1)$, and are confirmed by the automaton for $|n|,|j|\le12$.
\end{proof}

\begin{corollary}\label{cor:antipair}
For $h\ge2$, $\min_{(n,j)}\max\{k(n,j),k(n+h,j)\}=h-1$, attained on the row
$j=1$ at $n=-h/2-1$ when $h$ is even and on the row $j=0$ at $n=-(h+1)/2$ when
$h$ is odd.
\end{corollary}

\begin{proof}
On a row with minimum $k_j$ the two values are $k_j+2a$ and $k_j+2b$ with
$a+b\ge h-1$, $a,b\ge0$ integers, since the valley has width at most two and
each step adds $2$; hence the larger is at least $k_j+h-1$, and $k_j\ge1$ off
the rows $0,1$. On row $0$ the pair $(n,n+h)$ with $n\le-1\le0\le n+h$ has
values $-2n-2$ and $2(n+h)$, equal to $h-1$ exactly when $h$ is odd and
$n=-(h+1)/2$; on row $1$ the values are $-2n-3$ and $2(n+h)+1$, equal to $h-1$
exactly when $h$ is even and $n=-h/2-1$. Parity makes one of the two rows
attain $h-1$ for every $h$.
\end{proof}

\begin{proposition}[Antipodal lamps, $m$ even]\label{prop:antipodal}
Let $m$ be even. Then $(r_0r_1)^{m/2}$ is the half-turn about the apex, and
\[
 t_{n+m/2,\,j}\;=\;-\,t_{n,j}\qquad\text{for all }(n,j).
\]
Part (i) of Theorem~\ref{thm:cyl} follows by pairing $n$ with $n+m/2$. The
generic translations $t_{n,j}$ and $-t_{n+m/2,j}$, distinct for generic shapes,
are identified on the stratum, and the shortest such identification has length
$m+4$.
\end{proposition}

\begin{proof}
$r_0r_1$ is the rotation by $2\pi/m$ about the apex, so its $(m/2)$th power is
the rotation by $\pi$. If $h$ is a half-turn and $t$ the translation by $v$,
then $hth^{-1}$ is the translation by $-v$. As $t_{n+m/2,j}$ is the conjugate
of $t_{n,j}$ by $(r_0r_1)^{m/2}$, the relation follows; summing it over the
$m/2$ pairs $\{n,n+m/2\}$ of a level gives Theorem~\ref{thm:cyl}(i). A single
lamp has length $6+2k(n,j)$ by Corollary~\ref{cor:single}, so an identification
of this form first appears at $6+2\min\max\{k(n,j),k(n+m/2,j)\}$, the minimum
over pairs, which is $m/2-1$ by Corollary~\ref{cor:antipair} with $h=m/2$;
hence $m+4$.
\end{proof}

For even $m\ge4$ this precedes the length $2m+2$ at which new translations
appear, so the identifications are what the census sees first, and it falls
below the generic one; for odd $m$ there is no such pairing and the census
rises instead.

\begin{proposition}[Stratum translation census]\label{prop:cylcensus}
For $\tau$ generic on the stratum $\alpha=\pi/m$, the numbers of translations
of each length are
\[
\begin{array}{r|rrrrrrr}
 \text{length} & 6&8&10&12&14&16&18\\\hline
 \text{generic} & 6&6&42&96&350&1092&3684\\
 m=3 & 6&8&38&76&224&630&1456\\
 m=4 & 6&4&34&62&198&480&1394\\
 m=5 & 6&6&42&104&448&1274&4612\\
 m=6 & 6&6&40&78&278&750&2342\\
 m=7 & 6&6&42&96&350&1110&4204\\
 m=9,11 & 6&6&42&96&350&1092&3684
\end{array}
\]
For odd $m\ge3$ the two censuses agree up to length $2m+1$ and the stratum
acquires translations that are not images of generic ones at $2m+2$; for even
$m\ge4$ they first differ at $m+4$, by Proposition~\ref{prop:antipodal}.
\end{proposition}

\begin{proof}
By Theorem~\ref{thm:rot}(i) a word represents a translation on the stratum
exactly when $c_2=0$ and $m\mid c_1$. If $c_1=0$ the linear part is trivial at
every shape, so the word represents a generic translation as well. Otherwise
$|c_1|\ge m$, so by Theorem~\ref{thm:rot}(ii) the word has length at least
$2m$, with equality only for $(r_0r_1)^{\pm m}$, which is the identity on the
stratum. Hence the shortest translation of the stratum that is not the image
of a generic one has length $2m+2$. For odd $m$, identifications among generic
translations require the level relation, hence length at least
$4\lceil m/2\rceil+2=2m+4$ by the arc count above, while for even $m$
Proposition~\ref{prop:antipodal} places the first identification at $m+4$.
Shortening is excluded below $2m+2$ as well: if $g\neq1$ has stratum length
$L\le2m$, then a stratum-geodesic word for $g$ has $c_1=0$, since otherwise
$|c_1|\ge m$ forces length at least $2m$ with equality only at the identity;
that word is then a word for a generic translation, so its generic length is at
most $L$, and the reverse inequality is trivial. For odd $m$ none of the three
effects is available below $2m+2$ and the censuses agree length by length,
while at $2m+2$ the first new translations appear. The displayed values are
computed by exact rational breadth-first search in the stratum group, which at
any single shape bounds the counts from below; they are confirmed at two leg
samples through length $18$ for $m\le9$, and at one sample for $m=11$. The
row $m=9,11$ reproduces the generic census exactly, as the onset statement
requires, $2m+2$ exceeding $18$ for both.
\end{proof}

For $m=5$ the stratum has \emph{more} translations than a generic shape at
the length where they first differ, the new translations outnumbering the
identifications; for $m=3$, where the levels are short, the identifications
dominate and the census falls below the generic one already at length $10$.
Beyond length $14$ we do not report values: the configuration enumeration on
$X_m$ is not known to be exhaustive there, and for $m=9$ it undercounts, the
correct value at length $18$ being the generic $3684$ by the argument of the
proposition.

\begin{remark}
Rational leg samples that look generic need not be. For $m=3$ the legs $(1,2)$
give the right triangle with angles $\pi/6,\pi/3,\pi/2$, and for $m=5$ the legs
$(2,3)$ produce $22$ coincidences at length $14$ that three other samples do
not. Every computation reported here was repeated at three samples.
\end{remark}

\begin{remark}
These lines lie outside the \emph{typical} stratum of \cite{IP}, where the
angles are $\Q$-linearly independent and the shortest non-generic relations have
length $18$. On the lines $\alpha=\pi/m$ the minimum drops as low as $12$,
attained by the relators of the three $m=4$ relations displayed above.
\end{remark}

\section*{Declarations}
\noindent\textbf{Funding.} No funding was received for this work.
\smallskip\\
\noindent\textbf{Competing interests.} The author declares no competing interests.
\smallskip\\
\noindent\textbf{Formalization.} The arithmetic half of the lower bound of
Theorem~\ref{thm:metric} is formalized in Lean~4 in
\texttt{lean/TriangleFlowMetric.lean} of the cited repository, with no
imports, no \texttt{sorry}, and no axioms beyond \texttt{propext} and
\texttt{Quot.sound}. The correspondence is
\texttt{flowNorm\_le\_traversals} and \texttt{traversals\_parity} for the two
per-edge facts, \texttt{connector\_two\_le} for the doubling of an edge
carrying no flow, \texttt{steps\_ge} for the summation, and
\texttt{lower\_bound} for the bound itself, the graph-theoretic half entering
as its hypothesis \texttt{hst}. The remaining parts of the paper are not
formalized.
\smallskip\\
\noindent\textbf{Data availability.} All code and verification data supporting
the computational statements are openly available in the repository cited on the
title page and archived on Zenodo (concept DOI
\texttt{10.5281/zenodo.20370090}).
\smallskip\\
\noindent\textbf{Use of AI tools.} Large language models (Anthropic Claude) were
used in the preparation of this manuscript, both in drafting and in computation.
Every computational claim is certified by the scripts in the cited repository,
independently of any model output. The author takes full responsibility for the
content.

\begin{thebibliography}{9}
\bibitem{IP} S.~Isola and R.~Piergallini, \emph{On the generic triangle group},
arXiv:1405.1881 [math.MG], 2014--2015.
\bibitem{Bonfioli1} V.~Bonfioli, \emph{The universal right-triangle reflection
sequence: a word-length metric, effective universality, and the lamplighter
structure of A396406}, preprint (2026),
\texttt{doi:10.5281/zenodo.20370090}.
\end{thebibliography}

\end{document}
